% ---------------------------------------------------------------------
\section{Additional results}\label{sec:results-supp}
% ---------------------------------------------------------------------
% SKELETON. Tables and figures are read from ../../analysis/ and appear
% automatically once the analysis scripts have run on the real results.

\TODO{Skeleton: fill the result placeholders from the analysis outputs, use the past tense and write the interpretation yourself}.

\subsection{Learning curves and full result tables}\label{app:res-tables}

\Fig{fig:curves} shows the training curves of the learned methods, and \Tab{tab:t2} lists the zero-shot results behind \Fig{M-fig:zeroshot} for every station size.
\Tab{tab:t3} gives the constraint-satisfaction metrics behind \Fig{M-fig:constraints}, and \Tab{tab:t5} gives the cost of the guarantee behind \Fig{M-fig:cost}.
Across seeds, the training curves of \method{} \RES{description of convergence and seed spread}.

\begin{figure}[t]
\centering
\resultfig{../../analysis/figures/figS1_training_curves.pdf}{Episode profit and constraint metrics versus environment steps for the learned methods, mean and spread over seeds.}
\caption{Training curves of the learned methods.}
\label{fig:curves}
\end{figure}

\resulttable{../../analysis/tables/T2.tex}{\placeholdertable{tab:t2}{Zero-shot transfer across station sizes $N$}{T2.tex}}
\resulttable{../../analysis/tables/T3.tex}{\placeholdertable{tab:t3}{Constraint satisfaction per method over all test-grid configurations}{T3.tex}}
\resulttable{../../analysis/tables/T5.tex}{\placeholdertable{tab:t5}{Cost of the guarantee per station size $N$}{T5.tex}}

\subsection{Ablation study}\label{app:res-ablation}

\Fig{fig:ablation} and \Tab{tab:t4} compare attention and MLP policies with and without the safety layer, together with the further ablations of \method{}.
Without admission control, infeasible steps and shortfall change by \RES{change at the training sizes and at $N=200$}.
Setting $\beta=0$ changes the projection distance by \RES{value} and the profit by \RES{value}, and the curriculum and the small network give \RES{gaps relative to the full method}.

\begin{figure}[t]
\centering
\resultfig{../../analysis/figures/figS2_ablation.pdf}{Optimality gap and constraint metrics of the ablation variants.}
\caption{Ablation study.}
\label{fig:ablation}
\end{figure}

\resulttable{../../analysis/tables/T4.tex}{\placeholdertable{tab:t4}{Ablation of architecture, constraint handling and training choices}{T4.tex}}

\subsection{Conformal departure bound}\label{app:res-conformal}

\Fig{fig:conformal} and \Tab{tab:t7} show coverage, service and profit of \method{}-CP against the miscoverage level $\alpha$.
The episode-averaged coverage is \RES{coverage for each $\alpha$} against the nominal levels $1-\alpha$.
Compared with known departure steps, the profit falls by \RES{profit loss for each $\alpha$}, whereas the uncalibrated point prediction gives \RES{coverage and shortfall of the point prediction}.

\begin{figure}[t]
\centering
\resultfig{../../analysis/figures/figS3_conformal_alpha.pdf}{Coverage, energy not served and profit loss versus $\alpha$.}
\caption{Trade-off of the conformal departure bound.}
\label{fig:conformal}
\end{figure}

\resulttable{../../analysis/tables/T7.tex}{\placeholdertable{tab:t7}{Conformal departure bound: miscoverage level $\alpha$ versus profit and departure shortfall}{T7.tex}}

\subsection{Computation time}\label{app:res-time}

\Fig{fig:inference} and \Tab{tab:t6} report the inference time per control step of the policy and the safety layer, together with the training time per run.
At $N=10$ and $N=200$, one step takes \RES{mean inference time in ms}, of which the safety layer takes \RES{share}.
Arrival steps with admission control take \RES{time of the admission LP at $N=200$}, and one oracle LP takes \RES{solve time}.

\begin{figure}[t]
\centering
\resultfig{../../analysis/figures/figS4_inference_time.pdf}{Inference time per control step versus $N$ for policy and safety layer.}
\caption{Inference time per control step.}
\label{fig:inference}
\end{figure}

\resulttable{../../analysis/tables/T6.tex}{\placeholdertable{tab:t6}{Computation: training time per run and inference time per control step}{T6.tex}}

\subsection{Optimality gap with respect to both optima and sensitivity studies}\label{app:res-sensitivity}

\Tab{tab:ts1} reports the optimality gap with respect to $J^{\sm{LP,ceil}}$ and to $J^{\sm{LP}}$ for every station size, and the difference between the two optima is \RES{mean difference $J^{\sm{LP,ceil}}-J^{\sm{LP}}$ per scenario}.
In the sensitivity studies of \Tab{tab:sensitivity}, a desired state of charge of 0.8 gives \RES{gap and shortfall}, and a step of 5 min gives \RES{gap and shortfall}.

\resulttable{../../analysis/tables/TS1.tex}{\placeholdertable{tab:ts1}{Optimality gap with respect to $J^{\sm{LP,ceil}}$ and $J^{\sm{LP}}$ per station size $N$}{TS1.tex}}

\begin{table}[t]
\centering
\caption{Sensitivity studies at $N=20$ and $\rho_{\sm{tr}}=4$ kW (all three profiles).}
\label{tab:sensitivity}
\footnotesize
\begin{tabular}{@{}lcccc@{}}
\toprule
Setting & Gap [\%] & Unserved [\%] & Violation [kWh] & Infeasible \\
\midrule
Main setting & \RES{} & \RES{} & \RES{} & \RES{} \\
Desired state of charge 0.8 & \RES{} & \RES{} & \RES{} & \RES{} \\
$\Delta t=5$ min & \RES{} & \RES{} & \RES{} & \RES{} \\
Penalty weights 1 and 100 \texteuro/kWh & \RES{} & \RES{} & \RES{} & \RES{} \\
\bottomrule
\end{tabular}
\end{table}

\subsection{Pipeline and example episode}\label{app:res-example}

\Fig{fig:pipeline} summarises the experimental pipeline, and \Fig{fig:exampleday} traces one test episode of \method{} with the aggregate power against the forward and reverse limits, the band between $\sum_il_i$ and $\sum_iu_i$, and the price.
In this episode, \method{} \RES{description of how charging and V2G discharging follow the price and the limits}.

\begin{figure}[t]
\centering
\figchain{0.9\linewidth}{../../analysis/figures/figS5_pipeline.pdf}{../../analysis/figures_conceptual/figS5_pipeline.pdf}{figures/figS5_pipeline.pdf}
\caption{Experimental pipeline. (1) The learned methods train on mixed station sizes, arrival profiles and ratings. (2) A zero-shot test grid has 36 configurations plus the literature configuration, with sizes seen in training shaded. (3) All methods run on identical replays. (4) Metrics are computed per episode, and the perfect-information LP provides the reference optimum of every replay.}
\label{fig:pipeline}
\end{figure}

\begin{figure}[t]
\centering
\resultfig{../../analysis/figures/fig5_example_day.pdf}{Aggregate power, $C(t)$, $-\Crev(t)$, the band between $\sum_il_i$ and $\sum_iu_i$, and the price over one episode; fallback steps marked.}
\caption{Example episode of \method{} (workplace profile, $N=20$, tight rating).}
\label{fig:exampleday}
\end{figure}
